\chapter{2006年安徽卷（理）}
\section{单选题}
\begin{problem}
复数 \(\frac{1+\sqrt{3}\i}{\sqrt{3}-\i}\) 等于（\(\quad\)）

\choices
    {\(\i\)}
    {\(-\i\)}
    {\(\sqrt{3}+\i\)}
    {\(\sqrt{3}-\i\)}
\end{problem}

\begin{answer}
A
\end{answer}

\begin{solution}
将分子、分母同乘分母的共轭复数 \(\sqrt{3}+\i\)，得
\[
\frac{1+\sqrt{3}\i}{\sqrt{3}-\i}
=\frac{(1+\sqrt{3}\i)(\sqrt{3}+\i)}{(\sqrt{3}-\i)(\sqrt{3}+\i)}
=\frac{\sqrt{3}+\i+3\i-\sqrt{3}}{3+1}
=\frac{4\i}{4}=\i
\]
因此选 A.
\end{solution}

\begin{problem}
设集合 \(A = \{x \mid |x - 2| \leqslant 2, x \in \R\}\)，\(B = \{y \mid y = -x^2, -1 \leqslant x \leqslant 2\}\)，则 \(\complement_{\R}(A \cap B)\) 等于（\(\quad\)）

\choices
    {\(\R\)}
    {\(\{x\mid x\in \R,x\neq 0\}\)}
    {\(\{0\}\)}
    {\(\varnothing\)}
\end{problem}

\begin{answer}
B
\end{answer}

\begin{solution}
由 \(|x-2|\leqslant2\)，得 \(0\leqslant x\leqslant4\)，所以 \(A=[0,4]\).
又因为 \(y=-x^2\)，且 \(-1\leqslant x\leqslant2\)，函数值的最大值为 \(0\)，最小值为 \(-4\)，所以 \(B=[-4,0]\).
从而 \(A\cap B=\{0\}\)，其在 \(\R\) 中的补集为 \(\R\mathbin{\backslash}\{0\}\)，因此选 B.
\end{solution}

\begin{problem}
若抛物线 \(y^{2} = 2px\) 的焦点与椭圆 \(\frac{x^2}{6} + \frac{y^2}{2} = 1\) 的右焦点重合，则 \(p\) 的值为（\(\quad\)）

\choices
    {\(-2\)}
    {\(2\)}
    {\(-4\)}
    {\(4\)}
\end{problem}

\begin{answer}
D
\end{answer}

\begin{solution}
椭圆的半长轴平方、半短轴平方分别为 \(6\)、\(2\)，故
\[
c^2=6-2=4
\]
其右焦点为 \((2,0)\). 抛物线 \(y^2=2px\) 的焦点为 \((p/2,0)\)，由焦点重合得
\[
\frac p2=2
\]
所以 \(p=4\)，因此选 D.
\end{solution}

\begin{problem}
设 \(a\)，\(b \in \R\)，已知命题 \(p\)：\(a = b\)，命题 \(q\)：\(\left(\frac{a + b}{2}\right)^2 \leqslant \frac{a^2 + b^2}{2}\)，则 \(p\) 是 \(q\) 成立的（\(\quad\)）

\choices
    {必要不充分条件}
    {充分不必要条件}
    {充分必要条件}
    {既不充分也不必要条件}
\end{problem}

\begin{answer}
B
\end{answer}

\begin{solution}
将命题 \(q\) 两边相减并通分，得到
\[
\frac{a^2+b^2}{2}-\left(\frac{a+b}{2}\right)^2
=\frac{(a-b)^2}{4}\geqslant0
\]
所以对任意实数 \(a\)，\(b\)，命题 \(q\) 恒成立. 若 \(p\) 成立，即 \(a=b\)，当然有 \(q\) 成立；但 \(q\) 成立不要求 \(a=b\)，例如 \(a=1\)，\(b=0\) 时仍成立. 因此 \(p\) 是 \(q\) 成立的充分不必要条件，选 B.
\end{solution}

\begin{problem}
函数 \(y = \begin{cases}
2x, & x \geqslant 0,\\
-x^2, & x < 0,
\end{cases}\) 的反函数是（\(\quad\)）

\choices
    {\(y = \begin{cases}
\frac{x}{2}, & x \geqslant 0,\\
\sqrt{-x}, & x < 0.
\end{cases}\)}
    {\(y = \begin{cases}
2x, & x \geqslant 0,\\
\sqrt{-x}, & x < 0.
\end{cases}\)}
    {\(y = \begin{cases}
\frac{x}{2}, & x \geqslant 0,\\
-\sqrt{-x}, & x < 0.
\end{cases}\)}
    {\(y = \begin{cases}
2x, & x \geqslant 0,\\
-\sqrt{-x}, & x < 0.
\end{cases}\)}
\end{problem}

\begin{answer}
C
\end{answer}

\begin{solution}
原函数在 \(x\geqslant0\) 时的值域为 \([0,+\infty)\)，在 \(x<0\) 时的值域为 \(( -\infty,0)\).
对反函数设自变量为 \(x\). 当 \(x\geqslant0\) 时，原式 \(x=2y\)，故 \(y=x/2\)；当 \(x<0\) 时，原式 \(x=-y^2\)，且原函数这一段要求 \(y<0\)，故 \(y=-\sqrt{-x}\).
所以反函数为
\[
y=\begin{cases}
\dfrac{x}{2},&x\geqslant0,\\
-\sqrt{-x},&x<0,
\end{cases}
\]
选 C.
\end{solution}

\begin{problem}
将函数 \(y = \sin \omega x\)（\(\omega > 0\)）的图象按向量 \(\bs{a} = \left(-\frac{\pi}{6}, 0\right)\) 平移，平移后的图象如图所示，则平移后的图象所对应的函数解析式是（\(\quad\)）

\bitmapfigure[max width=0.42\linewidth]{img/2006/anhui_science/q6_fig1.png}

\choices
    {\(y = \sin \left(x + \frac{\pi}{6}\right)\)}
    {\(y = \sin \left(x - \frac{\pi}{6}\right)\)}
    {\(y = \sin \left(2x + \frac{\pi}{3}\right)\)}
    {\(y = \sin \left(2x - \frac{\pi}{3}\right)\)}
\end{problem}

\begin{answer}
C
\end{answer}

\begin{solution}
图象向左平移 \(\pi/6\) 后，原函数的自变量应替换为 \(x+\pi/6\)，故平移后的函数形式为
\[
y=\sin\omega\left(x+\frac{\pi}{6}\right)
\]
由图象相邻最高点与最低点的横坐标相差 \(\pi/2\)，这是半个周期，所以周期为 \(\pi\)，从而 \(2\pi/\omega=\pi\)，得 \(\omega=2\). 因此
\[
y=\sin\left(2x+\frac{\pi}{3}\right)
\]
选 C.
\end{solution}

\begin{problem}
若曲线 \(y = x^4\) 的一条切线 \(l\) 与直线 \(x + 4y - 8 = 0\) 垂直，则 \(l\) 的方程为（\(\quad\)）

\choices
    {\(4x - y - 3 = 0\)}
    {\(x + 4y - 5 = 0\)}
    {\(4x - y + 3 = 0\)}
    {\(x + 4y + 3 = 0\)}
\end{problem}

\begin{answer}
A
\end{answer}

\begin{solution}
直线 \(x+4y-8=0\) 的斜率为 \(-1/4\)，所以与它垂直的切线斜率为 \(4\). 曲线 \(y=x^4\) 在点 \((x,x^4)\) 处的切线斜率为
\[
y'=4x^3
\]
令 \(4x^3=4\)，得 \(x=1\)，切点为 \((1,1)\). 切线方程为
\[
y-1=4(x-1)
\]
即 \(4x-y-3=0\)，选 A.
\end{solution}

\begin{problem}
设 \(a > 0\)，对于函数 \(f(x) = \frac{\sin x + a}{\sin x}\)（\(0 < x < \pi\)），下列结论正确的是（\(\quad\)）

\choices
    {有最大值而无最小值}
    {有最小值而无最大值}
    {有最大值且有最小值}
    {既无最大值又无最小值}
\end{problem}

\begin{answer}
B
\end{answer}

\begin{solution}
因为 \(0<x<\pi\)，所以 \(0<\sin x\leqslant1\)，且 \(\sin x\) 在 \(x=\pi/2\) 时取最大值. 函数可写为
\[
f(x)=1+\frac{a}{\sin x}
\]
由于 \(a>0\)，有
\[
f(x)\geqslant1+a
\]
等号在 \(x=\pi/2\) 时成立，所以函数有最小值 \(1+a\). 当 \(x\to0^+\) 或 \(x\to\pi^-\) 时，\(\sin x\to0^+\)，故 \(f(x)\to+\infty\)，函数没有最大值. 因此选 B.
\end{solution}

\begin{problem}
表面积为 \(2\sqrt{3}\) 的正八面体的各个顶点都在同一球面上，则此球的体积为（\(\quad\)）

\choices
    {\(\frac{\sqrt{2}}{3}\pi\)}
    {\(\frac{1}{3}\pi\)}
    {\(\frac{2}{3}\pi\)}
    {\(\frac{2\sqrt{2}}{3}\pi\)}
\end{problem}

\begin{answer}
A
\end{answer}

\begin{solution}
设正八面体的棱长为 \(s\). 正八面体有 \(8\) 个全等的正三角形面，故其表面积为
\[
8\cdot\frac{\sqrt{3}}{4}s^2=2\sqrt{3}s^2
\]
由题设表面积为 \(2\sqrt{3}\)，得 \(s=1\). 正八面体的中心到任一顶点的距离为外接球半径
\[
R=\frac{s}{\sqrt{2}}=\frac{1}{\sqrt{2}}
\]
所以球的体积为
\[
V=\frac{4}{3}\pi R^3=\frac{4}{3}\pi\left(\frac{1}{\sqrt{2}}\right)^3=\frac{\sqrt{2}}{3}\pi
\]
因此选 A.
\end{solution}

\begin{problem}
如果实数 \(x\)，\(y\) 满足条件 \(\begin{cases}
x - y + 1 \geqslant 0,\\
y + 1 \geqslant 0,\\
x + y + 1 \leqslant 0,
\end{cases}\) 那么 \(2x - y\) 的最大值为（\(\quad\)）

\choices
    {\(2\)}
    {\(1\)}
    {\(-2\)}
    {\(-3\)}
\end{problem}

\begin{answer}
B
\end{answer}

\begin{solution}
三个不等式分别给出
\(y\leqslant x+1\)，\(y\geqslant-1\)，\(y\leqslant-x-1\).
可行域是由三条边界直线围成的三角形，其顶点为
\((-2,-1)\)，\((-1,0)\)，\((0,-1)\).
目标函数 \(2x-y\) 在这些顶点处的值依次为
\(-3\)，\(-2\)，\(1\).
线性目标函数在多边形上的最大值在顶点取得，所以最大值为 \(1\)，选 B.
\end{solution}

\begin{problem}
如果 \(\triangle A_{1}B_{1}C_{1}\) 的三个内角的余弦值分别等于 \(\triangle A_{2}B_{2}C_{2}\) 的三个内角的正弦值，则（\(\quad\)）

\choices
    {\(\triangle A_{1}B_{1}C_{1}\) 和 \(\triangle A_{2}B_{2}C_{2}\) 都是锐角三角形}
    {\(\triangle A_{1}B_{1}C_{1}\) 和 \(\triangle A_{2}B_{2}C_{2}\) 都是钝角三角形}
    {\(\triangle A_{1}B_{1}C_{1}\) 是钝角三角形，\(\triangle A_{2}B_{2}C_{2}\) 是锐角三角形}
    {\(\triangle A_{1}B_{1}C_{1}\) 是锐角三角形，\(\triangle A_{2}B_{2}C_{2}\) 是钝角三角形}
\end{problem}

\begin{answer}
D
\end{answer}

\begin{solution}
三角形 \(A_1B_1C_1\) 的三个内角余弦都等于正数，因此其三个内角都小于 \(\pi/2\)，所以 \(\triangle A_1B_1C_1\) 是锐角三角形.

若 \(\triangle A_2B_2C_2\) 也是锐角三角形，则各角都在 \((0,\pi/2)\) 内. 由
\[
\sin A_{2}=\cos A_{1}=\sin\left(\frac{\pi}{2}-A_{1}\right)
\]
以及正弦函数在 \((0,\pi/2)\) 上的单调性，得 \(A_2=\pi/2-A_1\). 对 \(B_2\) 同样有
\[
B_2=\frac{\pi}{2}-B_1
\]
同样由 \(\sin C_2=\cos C_1=\sin\left(\frac{\pi}{2}-C_1\right)\) 及正弦函数在 \((0,\pi/2)\) 上的单调性，得
\[
C_2=\frac{\pi}{2}-C_1
\]
于是 \(A_2+B_2+C_2=\pi/2\)，与三角形内角和为 \(\pi\) 矛盾. \(\triangle A_2B_2C_2\) 也不可能是直角三角形，因为若某个角为 \(\pi/2\)，其正弦为 \(1\)，则对应的 \(A_1\) 角余弦为 \(1\)，这要求 \(A_1=0\)，不可能.

因此 \(\triangle A_2B_2C_2\) 是钝角三角形，选 D.
\end{solution}

\begin{problem}
在正方体上任选 \(3\) 个顶点连成三角形，则所得的三角形是直角非等腰三角形的概率为（\(\quad\)）

\choices
    {\(\frac{1}{7}\)}
    {\(\frac{2}{7}\)}
    {\(\frac{3}{7}\)}
    {\(\frac{4}{7}\)}
\end{problem}

\begin{answer}
C
\end{answer}

\begin{solution}
从正方体的 \(8\) 个顶点中任取 \(3\) 个，样本总数为
\[
\mathrm C_{8}^{3}=56
\]
固定一个顶点作为直角顶点. 从该顶点出发的三个棱方向两两垂直. 若另外两个顶点分别在从该顶点出发的两条棱上，则有 \(3\) 种选法，但所得两条直角边等长；若一个顶点在一条棱的另一端，另一个顶点在不含该棱的那个面对角线的对角端，则也有 \(3\) 种选法，且两条直角边长度分别为 \(1\) 和 \(\sqrt{2}\)，所得三角形为非等腰直角三角形.

每个直角三角形的直角顶点唯一，故符合条件的三角形数为 \(8\times3=24\). 所求概率为
\[
\frac{24}{56}=\frac{3}{7}
\]
因此选 C.
\end{solution}

\section{填空题}
\begin{problem}
设常数 \(a > 0\)，\(\left(a x^2 + \frac{1}{\sqrt{x}}\right)^4\) 展开式中 \(x^3\) 的系数为 \(\frac{3}{2}\)，则 \(\displaystyle\lim_{n \to +\infty} (a + a^2 + \cdots + a^n) = \)\fillinblank{}.
\end{problem}

\begin{answer}
1
\end{answer}

\begin{solution}
在展开式中，若取 \(k\) 个 \(ax^2\)，取 \(4-k\) 个 \(1/\sqrt{x}\)，则该项中 \(x\) 的指数为
\[
2k-\frac{4-k}{2}=\frac{5k-4}{2}
\]
令其等于 \(3\)，得 \(k=2\). 因此 \(x^3\) 的系数为
\[
\mathrm C_{4}^{2}a^2=6a^2=\frac{3}{2}
\]
从而 \(a^2=1/4\). 由 \(a>0\)，得 \(a=1/2\). 所以
\[
\lim_{n\to+\infty}(a+a^2+\cdots+a^n)
=\lim_{n\to+\infty}\sum_{k=1}^{n}\left(\frac12\right)^k
=\frac{1/2}{1-1/2}=1
\]
\end{solution}

\begin{problem}
在平行四边形 \(ABCD\) 中，\(\overrightarrow{AB} = \bs{a}\)，\(\overrightarrow{AD} = \bs{b}\)，\(\overrightarrow{AN} = 3\overrightarrow{NC}\)，\(M\) 为 \(BC\) 的中点，则 \(\overrightarrow{MN} =\)\fillinblank{}.

（用 \(\bs{a}\)，\(\bs{b}\) 表示）
\end{problem}

\begin{answer}
\(-\frac{1}{4}\bs{a}+\frac{1}{4}\bs{b}\)
\end{answer}

\begin{solution}
取 \(A\) 为原点，则
\(\overrightarrow{AB}=\bs{a}\)，\(\overrightarrow{AD}=\bs{b}\)，\(\overrightarrow{AC}=\bs{a}+\bs{b}\).
由 \(\overrightarrow{AN}=3\overrightarrow{NC}\)，有
\[
\overrightarrow{AC}=\overrightarrow{AN}+\overrightarrow{NC}
=4\overrightarrow{NC}
\]
所以
\[
\overrightarrow{AN}=\frac34(\bs{a}+\bs{b})
\]
又因为 \(M\) 是 \(BC\) 的中点，
\[
\overrightarrow{AM}=\overrightarrow{AB}+\frac12\overrightarrow{BC}
=\bs{a}+\frac12\bs{b}
\]
因此
\[
\overrightarrow{MN}=\overrightarrow{AN}-\overrightarrow{AM}
=\frac34(\bs{a}+\bs{b})-\left(\bs{a}+\frac12\bs{b}\right)
=-\frac14\bs{a}+\frac14\bs{b}
\]
\end{solution}

\begin{problem}
函数 \(f(x)\) 对于任意实数 \(x\) 满足条件 \(f(x + 2) = \frac{1}{f(x)}\)，若 \(f(1) = -5\)，则 \(f(f(5)) = \)\fillinblank{}.
\end{problem}

\begin{answer}
\(-\frac{1}{5}\)
\end{answer}

\begin{solution}
由题设关系式可知函数值不为零，并且
\[
f(x+4)=\frac{1}{f(x+2)}=f(x)
\]
所以 \(f\) 以 \(4\) 为周期. 于是
\[
f(5)=f(1)=-5
\]
再由 \(f(1)=1/f(-1)\)，得
\[
f(-1)=-\frac15
\]
由于 \(-5=-1-4\)，利用周期性得到 \(f(-5)=f(-1)=-1/5\). 因此
\[
f(f(5))=f(-5)=-\frac15
\]
\end{solution}

\begin{problem}
多面体上，位于同一条棱两端的顶点称为相邻的. 如图，正方体的一个顶点 \(A\) 在平面 \(\alpha\) 内，其余顶点在 \(\alpha\) 的同侧，正方体上与顶点 \(A\) 相邻的三个顶点到 \(\alpha\) 的距离分别为 \(1\)，\(2\) 和 \(4\). \(P\) 是正方体的其余四个顶点中的一个，则 \(P\) 到平面 \(\alpha\) 的距离可能是：

\begin{circlelist}
\item \(3\)
\item \(4\)
\item \(5\)
\item \(6\)
\item \(7\)
\end{circlelist}

以上结论正确的是\fillinblank{}.（写出所有正确结论的编号）

\begingroup
\leftskip=0pt
\rightskip=0pt
\bitmapfigure[max width=0.42\linewidth]{img/2006/anhui_science/q16_fig1.png}
\endgroup
\end{problem}

\begin{answer}
1, 3, 4, 5
\end{answer}

\begin{solution}
以顶点 \(A\) 为原点，沿从 \(A\) 出发的三条棱取向量 \(\bs{u}\)，\(\bs{v}\)，\(\bs{w}\). 设带符号距离平面 \(\alpha\) 的有向距离函数为 \(d\)，则它是顶点坐标的仿射一次函数. 由题设
\(d(A)=0\)，\(d(A+\bs{u})=1\)，\(d(A+\bs{v})=2\)，\(d(A+\bs{w})=4\).
由于其余顶点均在平面 \(\alpha\) 的同侧，\(d\) 在这些顶点上取正值. 对任意实数 \(r\)，\(s\)，\(t\)，仿射一次函数满足
\[
d(A+r\bs{u}+s\bs{v}+t\bs{w})
=d(A)+r\bigl(d(A+\bs{u})-d(A)\bigr)
+s\bigl(d(A+\bs{v})-d(A)\bigr)
+t\bigl(d(A+\bs{w})-d(A)\bigr)
\]
因此在各棱向量系数为 \(0\) 或 \(1\) 的顶点上，有向距离就是相应邻接顶点距离之和. 于是另外四个顶点到平面的距离分别为
\[
d(A+\bs{u}+\bs{v})=1+2=3
\]
\[
d(A+\bs{u}+\bs{w})=1+4=5
\]
\[
d(A+\bs{v}+\bs{w})=2+4=6
\]
以及
\[
d(A+\bs{u}+\bs{v}+\bs{w})=1+2+4=7
\]
所以可能的距离为 \(3\)，\(5\)，\(6\)，\(7\)，对应选项编号为 \(1\)，\(3\)，\(4\)，\(5\).
\end{solution}

\section{解答题}
\begin{problem}
已知 \(\frac{3\pi}{4} < \alpha < \pi\)，\(\tan \alpha + \cot \alpha = -\frac{10}{3}\).

\begin{enumerate}
\item 求 \(\tan \alpha\) 的值；
\item 求 \(\frac{5\sin^{2}\frac{\alpha}{2}+8\sin\frac{\alpha}{2}\cos\frac{\alpha}{2}+11\cos^{2}\frac{\alpha}{2}-8}{\sqrt{2}\sin\left(\alpha-\frac{\pi}{2}\right)}\) 的值.
\end{enumerate}
\end{problem}

\begin{solution}
\begin{enumerate}
\item 令 \(t=\tan\alpha\). 由 \(\cot\alpha=1/t\)，题设给出
\[
t+\frac1t=-\frac{10}{3}
\]
因为 \(\tan\alpha\neq0\)，两边同乘 \(3t\)，得
\[
3t^2+10t+3=0
\]
即
\[
(3t+1)(t+3)=0
\]
所以 \(t=-1/3\) 或 \(t=-3\). 又因为 \(3\pi/4<\alpha<\pi\)，故 \(-1<\tan\alpha<0\)，舍去 \(-3\)，得到
\[
\tan\alpha=-\frac13
\]

\item 记 \(s=\sin(\alpha/2)\)，\(c=\cos(\alpha/2)\). 利用二倍角公式，分子为
\[
5s^2+8sc+11c^2-8
=\frac52(1-\cos\alpha)+4\sin\alpha+\frac{11}{2}(1+\cos\alpha)-8
=3\cos\alpha+4\sin\alpha
\]
分母为
\[
\sqrt2\sin\left(\alpha-\frac\pi2\right)=-\sqrt2\cos\alpha
\]
由于 \(\alpha\in(3\pi/4,\pi)\)，有 \(\cos\alpha\neq0\)，所以原式等于
\[
-\frac{3\cos\alpha+4\sin\alpha}{\sqrt2\cos\alpha}
=-\frac{3+4\tan\alpha}{\sqrt2}
=-\frac{3-4/3}{\sqrt2}
=-\frac{5\sqrt2}{6}
\]
\end{enumerate}
\end{solution}

\begin{problem}
在添加剂的搭配运用中，为了找到最佳的搭配方案，需要对各种不同的搭配方式作比较，在试制某种牙膏新品种时，需要选用两种不同的添加剂. 现有芳香度分别为 \(0\)，\(1\)，\(2\)，\(3\)，\(4\)，\(5\) 的六种添加剂可供选用. 根据实验设计学原理，通常首先要随机选取两种不同的添加剂进行搭配实验. 用 \(\xi\) 表示所选用的两种不同的添加剂的芳香度之和.

\begin{enumerate}
\item 写出 \(\xi\) 的分布列；（以列表的形式给出结论，不必写计算过程）
\item 求 \(\xi\) 的数学期望 \(E(\xi)\).（要求写出计算过程或说明理由）
\end{enumerate}
\end{problem}

\begin{solution}
\begin{enumerate}
\item 从 \(0\)，\(1\)，\(2\)，\(3\)，\(4\)，\(5\) 中任选两个不同的数，共有 \(\mathrm C_{6}^{2}=15\) 种等可能的选法. 各和的组合数如下，因此分布列为
\begin{center}
\begin{tabular}{c|ccccccccc}
\(\xi\) & \(1\) & \(2\) & \(3\) & \(4\) & \(5\) & \(6\) & \(7\) & \(8\) & \(9\) \\
\hline
\(P\) & \(\frac{1}{15}\) & \(\frac{1}{15}\) & \(\frac{2}{15}\) & \(\frac{2}{15}\) & \(\frac{3}{15}\) & \(\frac{2}{15}\) & \(\frac{2}{15}\) & \(\frac{1}{15}\) & \(\frac{1}{15}\)
\end{tabular}
\end{center}
\item 每个数在所有无序二元组中出现的次数相同，且任取一组的两个数的和的期望等于两个位置的平均数之和. 因此
\[
E(\xi)=2\cdot\frac{0+1+2+3+4+5}{6}=2\cdot\frac{15}{6}=5
\]
\end{enumerate}
\end{solution}

\begin{problem}
如图，\(P\) 是边长为 \(1\) 的正六边形 \(ABCDEF\) 所在平面外一点，\(PA = 1\)，\(P\) 在平面 \(ABC\) 内的射影为 \(BF\) 的中点 \(O\).

\begin{enumerate}
\item 证明：\(PA \perp BF\)；
\item 求面 \(APB\) 与面 \(DPB\) 所成二面角的大小.
\end{enumerate}

\bitmapfigure[max width=0.43\linewidth]{img/2006/anhui_science/q19_fig1.png}
\end{problem}

\begin{solution}
\begin{enumerate}
\item 取正六边形所在平面为 \(z=0\)，令 \(O\) 为原点，并取坐标
\(A=\left(\frac12,0,0\right)\)，\(B=\left(0,\frac{\sqrt3}{2},0\right)\)，\(F=\left(0,-\frac{\sqrt3}{2},0\right)\)，\(D=\left(-\frac32,0,0\right)\).
因为 \(O\) 是点 \(P\) 在该平面内的射影，所以设 \(P=(0,0,h)\). 由 \(PA=1\) 得
\[
\left(\frac12\right)^2+h^2=1
\]
从而 \(h=\sqrt3/2\)，可取 \(P=(0,0,\sqrt3/2)\). 于是
\(\overrightarrow{PA}=\left(\frac12,0,-\frac{\sqrt3}{2}\right)\)，\(\overrightarrow{BF}=\left(0,-\sqrt3,0\right)\).
两向量的数量积为 \(0\)，故 \(PA\perp BF\).

\item 设二面角为 \(\theta\). 两平面均以 \(PB\) 为交线，取
\(\bs{u}=\overrightarrow{PA}=\left(\frac12,0,-\frac{\sqrt3}{2}\right)\)，\(\bs{v}=\overrightarrow{PB}=\left(0,\frac{\sqrt3}{2},-\frac{\sqrt3}{2}\right)\)
\[
\bs{w}=\overrightarrow{PD}=\left(-\frac32,0,-\frac{\sqrt3}{2}\right)
\]
平面 \(APB\) 和平面 \(DPB\) 的法向量可分别取
\[
\bs{n}_1=\bs{u}\times\bs{v}
=\left(\frac34,\frac{\sqrt3}{4},\frac{\sqrt3}{4}\right)
\]
\[
\bs{n}_2=\bs{w}\times\bs{v}
=\left(\frac34,-\frac{3\sqrt3}{4},-\frac{3\sqrt3}{4}\right)
\]
因此
\(\bs{n}_1\cdot\bs{n}_2=-\frac{9}{16}\)，\(|\bs{n}_1|=\frac{\sqrt{15}}4\)，\(|\bs{n}_2|=\frac{3\sqrt7}{4}\).
按半平面 \(PBA\) 与 \(PBD\) 的取向，所求二面角 \(\theta\in(0,\pi)\) 对应于上述两个法向量的夹角，
\[
\cos\theta=\frac{\bs{n}_1\cdot\bs{n}_2}{|\bs{n}_1||\bs{n}_2|}
=-\frac{3}{\sqrt{105}}
\]
又 \(0<\theta<\pi\)，故 \(\tan\theta=-4\sqrt6/3\)，从而
\[
\theta=\pi-\arctan\left(\frac{4\sqrt6}{3}\right)
\]
\end{enumerate}
\end{solution}

\begin{problem}
已知函数 \(f(x)\) 在 \(\R\) 上有定义，对任意实数 \(a > 0\) 和任意实数 \(x\)，都有 \(f(ax) = af(x)\).

\begin{enumerate}
\item 证明 \(f(0) = 0\)；
\item 证明 \(f(x) = \begin{cases}
kx, & x \geqslant 0,\\
hx, & x < 0,
\end{cases}\) 其中 \(k\) 和 \(h\) 均为常数；
\item 当第（2）问中的 \(k > 0\) 时，设 \(g(x) = \frac{1}{f(x)} + f(x)\)（\(x > 0\)），讨论 \(g(x)\) 在 \((0, +\infty)\) 内的单调性并求极值.
\end{enumerate}
\end{problem}

\begin{solution}
\begin{enumerate}
\item 取任意正数 \(a\)，令 \(x=0\)，则
\[
f(0)=f(a\cdot0)=af(0)
\]
取 \(a=2\)，得 \(f(0)=2f(0)\)，所以 \(f(0)=0\).

\item 令 \(k=f(1)\). 当 \(x>0\) 时，取题设中的 \(a=x\)，得到
\[
f(x)=f(x\cdot1)=xf(1)=kx
\]
令 \(h=-f(-1)\). 当 \(x<0\) 时取 \(a=-x>0\)，则
\[
f(x)=f((-x)(-1))=(-x)f(-1)=hx
\]
结合第（1）问，得到
\[
f(x)=\begin{cases}
kx,&x\geqslant0,\\
hx,&x<0,
\end{cases}
\]
其中 \(k\)，\(h\) 均为常数.

\item 当 \(k>0\) 且 \(x>0\) 时，\(f(x)=kx>0\)，所以
\[
g(x)=\frac{1}{kx}+kx
\]
求导得
\[
g'(x)=k-\frac{1}{kx^2}=\frac{k^2x^2-1}{kx^2}
\]
由于 \(kx>0\)，导数在 \(0<x<1/k\) 时为负，在 \(x>1/k\) 时为正. 因此 \(g(x)\) 在 \((0,1/k)\) 内单调递减，在 \((1/k,+\infty)\) 内单调递增，并在 \(x=1/k\) 处取得极小值
\[
g\left(\frac1k\right)=\frac{1}{k\cdot(1/k)}+k\cdot\frac1k=2
\]
当 \(x\to0^+\) 或 \(x\to+\infty\) 时，\(g(x)\to+\infty\)，所以没有最大值.
\end{enumerate}
\end{solution}

\begin{problem}
数列 \(\{a_{n}\}\) 的前 \(n\) 项和为 \(S_{n}\). 已知 \(a_{1} = \frac{1}{2}\)，\(S_{n} = n^{2}a_{n} - n(n - 1)\)，\(n = 1\)，\(2\)，\(\dots\).

\begin{enumerate}
\item 写出 \(S_{n}\) 与 \(S_{n-1}\) 的递推关系式（\(n \geqslant 2\)）. 并求 \(S_{n}\) 关于 \(n\) 的表达式；
\item 设 \(f_n(x) = \frac{S_n}{n} x^{n+1}\)，\(b_n = f_n'(p)\)（\(p \in \R\)），求数列 \(\{b_n\}\) 的前 \(n\) 项和 \(T_n\).
\end{enumerate}
\end{problem}

\begin{solution}
\begin{enumerate}
\item 由 \(a_n=S_n-S_{n-1}\) 及题设关系式，得
\[
S_n=n^2(S_n-S_{n-1})-n(n-1)
\]
即
\((n^2-1)S_n-n^2S_{n-1}=n(n-1)\)，\(n\geqslant2\).
等价地，递推关系为
\[
\frac{n+1}{n}S_n-\frac{n}{n-1}S_{n-1}=1
\]
令 \(u_n=S_n/n\)，则
\[
(n+1)u_n-nu_{n-1}=1
\]
又 \(u_1=S_1=1/2\). 用归纳法证明 \(u_n=n/(n+1)\)：当 \(n=1\) 时成立；若 \(u_{n-1}=(n-1)/n\)，则
\[
(n+1)u_n-n\cdot\frac{n-1}{n}=1
\]
从而 \(u_n=n/(n+1)\). 因此
\[
S_n=nu_n=\frac{n^2}{n+1}
\]

\item 由上式
\[
f_n(x)=\frac{S_n}{n}x^{n+1}=\frac{n}{n+1}x^{n+1}
\]
所以
\[
b_n=f_n'(p)=np^n
\]
于是 \(T_n=\sum_{k=1}^nkp^k\). 当 \(p\neq1\) 时，将等比数列求和式求导，或直接计算
\[
(1-p)T_n
=\bigl(p+2p^2+\cdots+np^n\bigr)
-\bigl(p^2+2p^3+\cdots+np^{n+1}\bigr)
=p+p^2+\cdots+p^n-np^{n+1}
\]
得到
\[
T_n=\frac{p-(n+1)p^{n+1}+np^{n+2}}{(1-p)^2}
\]
当 \(p=1\) 时，
\[
T_n=1+2+\cdots+n=\frac{n(n+1)}2
\]
故
\[
T_{n}=\begin{cases}
\displaystyle\frac{p-(n+1)p^{n+1}+np^{n+2}}{(1-p)^{2}}, & p\neq 1,\\
\displaystyle\frac{n(n+1)}{2}, & p=1.
\end{cases}
\]
\end{enumerate}
\end{solution}

\begin{problem}
如图，\(F\) 为双曲线 \(C\)：\(\frac{x^{2}}{a^{2}} - \frac{y^{2}}{b^{2}} = 1\)（\(a > 0\)，\(b > 0\)）的右焦点，\(P\) 为双曲线 \(C\) 右支上一点，且位于 \(x\) 轴上方，\(M\) 为左准线上一点，\(O\) 为坐标原点. 已知四边形 \(OFPM\) 为平行四边形，\(|PF| = \lambda |OF|\).

\begin{enumerate}
\item 写出双曲线 \(C\) 的离心率 \(e\) 与 \(\lambda\) 的关系式；
\item 当 \(\lambda = 1\) 时，经过焦点 \(F\) 且平行于 \(OP\) 的直线交双曲线于 \(A\)，\(B\) 点，若 \(|AB| = 12\)，求此时的双曲线方程.
\end{enumerate}

\bitmapfigure[max width=0.44\linewidth]{img/2006/anhui_science/q22_fig1.png}
\end{problem}

\begin{solution}
\begin{enumerate}
\item 设双曲线右焦点为 \(F=(c,0)\)，其中
\(c^2=a^2+b^2\)，\(e=\frac ca\).
左、右准线分别为 \(x=-a^2/c\) 和 \(x=a^2/c\). 由四边形 \(OFPM\) 为平行四边形，得
\(M=P-F\)，\(PM=OF=c\).
设线段 \(PM\) 与右准线交于 \(H\). 由题设位置关系，\(M\)，\(H\)，\(P\) 依次位于该水平线上，且
\(HM=\frac{a^2}{c}-\left(-\frac{a^2}{c}\right)=\frac{2a^2}{c}\)，\(PH=PM-HM=c-\frac{2a^2}{c}\).
双曲线的离心率定义给出 \(PF=e\,PH\). 结合 \(PF=\lambda OF=\lambda c\)，得到
\[
\lambda c=e\left(c-\frac{2a^2}{c}\right)
\]
代入 \(c=ae\)，并约去 \(a\)，得
\[
\lambda e=e^2-2
\]
即
\[
e^2-\lambda e-2=0
\]

\item 当 \(\lambda=1\) 时，由 \(e>1\) 得 \(e=2\)，所以
\(c=2a\)，\(b^2=c^2-a^2=3a^2\).
由平行四边形关系及右准线位置，点 \(P\) 的横坐标为
\[
x_P=c-\frac{a^2}{c}=\frac{3a}{2}
\]
代入双曲线方程并取上方的点，得
\(P=\left(\frac{3a}{2},\frac{\sqrt{15}a}{2}\right)\)，\(F=(2a,0)\).
过 \(F\) 且平行于 \(OP\) 的直线为
\[
y=\frac{\sqrt{15}}{3}(x-2a)
\]
将其代入 \(x^2/a^2-y^2/(3a^2)=1\)，得到交点横坐标满足
\[
4x^2+20ax-29a^2=0
\]
两根之差为 \(3\sqrt6a\)，而直线斜率为 \(\sqrt{15}/3\)，所以两交点间距离为
\[
|AB|=3\sqrt6a\sqrt{1+\frac{15}{9}}=12a
\]
由 \(|AB|=12\)，得 \(a=1\)，进而 \(b^2=3\). 故双曲线方程为
\[
\frac{x^2}{1^2}-\frac{y^2}{3}=1
\]
即
\[
x^2-\frac{y^2}{3}=1
\]
\end{enumerate}
\end{solution}
